% Part: computability
% Chapter: computability-theory
% Section: ce-closed-cup-cap

\documentclass[../../../include/open-logic-section]{subfiles}

\begin{document}

\olfileid{cmp}{thy}{clo}

\olsection[C.E. সেটের সংযোজন ও ছেদ]{গণনসাধ্যভাবে তালিকায়নযোগ্য
  সেট সংযোজন ও ছেদের অধীনে বদ্ধ}

নিচের উপপাদ্যটি গণনসাধ্যভাবে তালিকায়নযোগ্য সেটসমূহের কয়েকটি বদ্ধতা-বৈশিষ্ট্য দেয়।

\begin{thm}
ধরা যাক, $A$ ও $B$ গণনসাধ্যভাবে তালিকায়নযোগ্য। তাহলে $A \cap B$
ও $A \cup B$-ও তাই।
\end{thm}

\begin{proof}
\olref[eqc]{thm:ce-equiv} গণনসাধ্যভাবে তালিকায়নযোগ্য সেটের বিভিন্ন
চরিত্রায়ন ব্যবহার করার সুযোগ দেয়। উদাহরণ হিসেবে আমরা কয়েকটি ভিন্ন
প্রমাণ দেব।

প্রথম প্রমাণের জন্য ধরা যাক, গণনসাধ্য অপেক্ষক~$f$, $A$-কে তালিকায়িত
করে এবং গণনসাধ্য অপেক্ষক~$g$, $B$-কে তালিকায়িত করে। ধরি,
\begin{align*}
h(x) & = \umin{y}{(f(y) = x \lor g(y) = x)} \text{ এবং}\\
j(x) & = \umin{y}{(f((y)_0) = x \land g((y)_1) = x)}.
\end{align*}
তাহলে $A \cup B$ হলো $h$-এর সংজ্ঞাক্ষেত্র এবং $A \cap B$ হলো~$j$-এর
সংজ্ঞাক্ষেত্র।

\begin{explain}
গণনার ভাষায় যা ঘটছে তা হলো: $A$ ও $B$-কে তালিকায়িত করার প্রণালী
দেওয়া থাকলে, যেকোনো একটি তালিকায়নে~$x$ খুঁজে আমরা $x$, $A \cup B$-তে
আছে কি না অর্ধ-নির্ণয় করতে পারি; এবং একই সঙ্গে উভয় তালিকায়নে~$x$
খুঁজে আমরা $x$, $A \cap B$-তে আছে কি না অর্ধ-নির্ণয় করতে পারি।
\end{explain}

দ্বিতীয় প্রমাণের জন্য আবার ধরা যাক, $f$, $A$-কে এবং $g$, $B$-কে
তালিকায়িত করে। ধরি,
\[
k(x) = \begin{cases}
f(x/2) & \text{যদি $x$ জোড় হয়} \\
g((x-1)/2) & \text{যদি $x$ বিজোড় হয়।}
\end{cases}
\]
তাহলে $k$, $A \cup B$-কে তালিকায়িত করে; ধারণাটি হলো, $k$ শুধু $f$ ও
~$g$-এর দেওয়া তালিকায়ন দুটির মধ্যে পালা করে। $A \cap B$-কে
তালিকায়িত করা একটু কঠিন। $A \cap B$ শূন্য হলে সেটি তুচ্ছভাবেই
গণনসাধ্যভাবে তালিকায়নযোগ্য। অন্যথায়, $c$-কে $A \cap B$-এর যেকোনো
উপাদান ধরি এবং~$l$-কে সংজ্ঞায়িত করি:
\[
l(x) = \begin{cases}
f((x)_0) & \text{যদি $f((x)_0) = g((x)_1)$} \\
c & \text{অন্যথায়।}
\end{cases}
\]
গণনার ভাষায়, $l$, $f$ ও $g$-এর তালিকায়ন থেকে উপাদানের জোড়াগুলির
মধ্য দিয়ে যায় এবং পাওয়া প্রতিটি মিল নির্গত করে; অন্যথায় শুধু~$c$
নির্গত করে অপেক্ষা করতে থাকে।

শেষ প্রমাণের জন্য ধরা যাক, $A$ হলো আংশিক অপেক্ষক~$m(x)$-এর
\emph{সংজ্ঞাক্ষেত্র} এবং $B$ হলো আংশিক অপেক্ষক~$n(x)$-এর সংজ্ঞাক্ষেত্র।
তাহলে $A \cap B$ হলো আংশিক অপেক্ষক $m(x) + n(x)$-এর সংজ্ঞাক্ষেত্র।

\begin{explain}
গণনার ভাষায়, $A$ যদি সেই সব মানের সেট হয় যাদের জন্য~$m$ থামে এবং
$B$ সেই সব মানের সেট হয় যাদের জন্য~$n$ থামে, তবে $A \cap B$ হলো সেই
সব মানের সেট যাদের জন্য উভয় প্রণালীই থামে।
\end{explain}

$A \cup B$-কে থামা-মানের সেট হিসেবে প্রকাশ করা আরও কঠিন, কারণ $m$ ও
$n$-কে সমান্তরালে অনুকরণ করতে হয়। ধরা যাক, $d$ হলো~$m$-এর এবং $e$
হলো~$n$-এর একটি সূচক; অন্যভাবে বললে, $m = \cfind{d}$ এবং
$n = \cfind{e}$। তাহলে $A \cup B$ নিচের অপেক্ষকটির সংজ্ঞাক্ষেত্র:
\[
p(x) = \umin{y}{(T(d,x,y) \lor T(e,x,y))}.
\]
\begin{explain}
গণনার ভাষায়, নিবেশ~$x$-এ $p$, $m$-এর অথবা $n$-এর একটি থামা গণনা
খোঁজে এবং যেকোনো একটি পেলেই থামে।
\end{explain}
\end{proof}

\end{document}
